\chapter*{Notation and Convention Guide}
\phantomsection
\label{ch:notation-convention-guide}
\addcontentsline{toc}{chapter}{Notation and Convention Guide}

This guide records global notation conventions used repeatedly across the
monograph.  It is a navigation aid, not a replacement for local definitions:
when a chapter first uses a load-bearing symbol it must still declare the
object, its ambient space or category, and its status.  If a symbol has
different meanings in different frameworks, the chapter must state the local
meaning before use.  The entries below therefore function as a consistency
check for recurrent notation and as a quick map for readers moving between
volumes.

\section*{Global Conventions}
\addcontentsline{toc}{section}{Global Conventions}

\begin{longtable}{>{\raggedright\arraybackslash}p{0.18\textwidth}
                  >{\raggedright\arraybackslash}p{0.74\textwidth}}
\hline
\textbf{Symbol or convention} & \textbf{Meaning and scope}\\
\hline
\endhead
\(\mathbb M^D\) &
\(D\)-dimensional Minkowski spacetime with metric
\(\eta_{\mu\nu}=\operatorname{diag}(-,+,\ldots,+)\), unless a chapter states
otherwise.  The monograph does not use mostly-minus conventions.\\
\(\mathbb R^D\) &
Euclidean \(D\)-space with positive metric \(\delta_{\mu\nu}\).  Euclidean
correlation functions are not automatically probability moments; their status
is the one declared in the local reconstruction or regulator datum.\\
\(\dd x\), \(\dd^D x\) &
Lebesgue measure in the displayed coordinate chart.  On manifolds the measure
is written with the appropriate density, volume form, or regulator datum.\\
\(\ii\), \(\ee\), \(\id\) &
The imaginary unit, exponential, and identity map or identity operator.
Identity operators on a named space may also be written \(\mathbf1\).\\
\(\hbar\) &
Set to \(1\) unless explicitly displayed.  Semiclassical discussions restore
\(\hbar\) locally and state the scaling regime.\\
\(U(a)=\exp(\ii a^\mu P_\mu)\) &
Translation convention.  A covariant field satisfies
\([P_\mu,\widehat\Phi(x)]=\ii\partial_\mu\widehat\Phi(x)\) on the common
invariant domain when point-field notation denotes a smeared
operator-valued distribution.\\
\([A,B]_\pm\) &
Graded commutator when a parity has been declared; ordinary commutator when
all operators are even.  Fermionic operator anticommutation is a statement
about Hilbert-space operators, distinct from Grassmann path-integral
variables.\\
\hline
\end{longtable}

\section*{Hilbert-Space and Operator-Algebra Symbols}
\addcontentsline{toc}{section}{Hilbert-Space and Operator-Algebra Symbols}

\begin{longtable}{>{\raggedright\arraybackslash}p{0.18\textwidth}
                  >{\raggedright\arraybackslash}p{0.74\textwidth}}
\hline
\textbf{Symbol} & \textbf{Meaning and scope}\\
\hline
\endhead
\(\Hilb\) &
Hilbert space of a reconstructed or assumed quantum theory.  Domains of
unbounded operators are part of the data whenever they enter a proof.\\
\(\Omega\) or \(\vac\) &
Vacuum vector when the Hilbert-space representation has been constructed.
Before reconstruction, brackets denote the stated Euclidean, algebraic, or
regulated functional expectation.\\
\(\Obs\), \(\mathcal A\) &
Observable algebra.  In Wightman chapters this may denote the algebra
generated by smeared fields on a common domain; in AQFT chapters
\(\Obs(\mathcal O)\) denotes the local algebra assigned to a spacetime
region \(\mathcal O\).\\
\(\mathcal R(\mathcal O)\) &
Von Neumann local algebra in a Hilbert-space representation, usually the
weak closure of \(\Obs(\mathcal O)\).\\
\(\Obs_{\rm qloc}\) &
Quasilocal algebra obtained as the inductive \(C^*\)-completion of local
observable algebras when such a completion is part of the stated AQFT datum.\\
\(\widehat\Phi(f)\) &
Smeared field operator.  The test function \(f\), target representation, and
common invariant domain must be specified locally.  Point notation
\(\widehat\Phi(x)\) is kernel notation subordinate to smearing.\\
\(E(\Delta)\) &
Spectral projection of a self-adjoint operator or commuting family for the
Borel set \(\Delta\).  For energy-momentum spectra the relevant joint
spectral measure is declared in the chapter.\\
\(\Hilb_1\) &
One-particle subspace associated with an isolated mass shell or other
declared spectral component.  It is not defined by a perturbative pole alone.\\
\(\mathcal S\), \(\mathcal S'\) &
Schwartz test functions and tempered distributions.  Distributional
equalities mean equality after pairing with test functions.\\
\(\mathcal D\), \(\mathcal D'\) &
Compactly supported smooth test functions and distributions when local
support, wavefront set, or microlocal arguments require compact support.\\
\hline
\end{longtable}

\section*{Correlation, Scattering, and Spectral Symbols}
\addcontentsline{toc}{section}{Correlation, Scattering, and Spectral Symbols}

\begin{longtable}{>{\raggedright\arraybackslash}p{0.18\textwidth}
                  >{\raggedright\arraybackslash}p{0.74\textwidth}}
\hline
\textbf{Symbol} & \textbf{Meaning and scope}\\
\hline
\endhead
\(W_n\) &
Lorentzian Wightman distribution.  Its arguments are ordered, and its
analytic continuation properties are distributional boundary-value
statements, not pointwise functions unless the chapter proves a stronger
regularity statement.\\
\(S_n\) &
Euclidean Schwinger distribution or regulated Euclidean correlator, depending
on the declared datum.  The symbol is not a literal moment of a Borel measure
unless such a measure has been constructed.\\
\(G_n\), \(G_n^{\rm conn}\) &
Time-ordered Green function and its connected part, with Lorentzian or
Euclidean status stated locally.\\
\(\rho(\mu^2)\), \(\dd\rho(\mu^2)\) &
K\"all\'en--Lehmann spectral density or positive spectral measure for a
gauge-invariant Hilbert-space two-point function.  Gauge-variant fields do
not inherit this positivity statement.\\
\(Z_\phi^{\rm pole}\) &
Residue of an isolated one-particle pole in the K\"all\'en--Lehmann sense.  It is
kept distinct from a generating functional \(Z[J]\).\\
\(\mathcal S\), \(S_{\beta\alpha}\) &
Scattering operator or its matrix elements after Haag--Ruelle wave operators
and LSZ hypotheses have been stated.  Perturbative diagrammatics for
S-matrix elements are downstream of this nonperturbative construction.\\
\(\mathcal M(s,t)\) &
Invariant amplitude in a declared normalization and sheet.  A chapter must
state whether the object is connected, amputated, physical-sheet,
second-sheet, partial-wave, or analytically continued.\\
\(s,t,u\) &
Mandelstam invariants, with metric convention fixed by the local chapter.
Analytic continuation in these variables is always tied to a sheet
convention.\\
\hline
\end{longtable}

\section*{Path-Integral, Effective-Action, and RG Symbols}
\addcontentsline{toc}{section}{Path-Integral, Effective-Action, and RG Symbols}

\begin{longtable}{>{\raggedright\arraybackslash}p{0.18\textwidth}
                  >{\raggedright\arraybackslash}p{0.74\textwidth}}
\hline
\textbf{Symbol} & \textbf{Meaning and scope}\\
\hline
\endhead
\([Dq]\), \([D\phi]\), \(D_\Lambda\phi\) &
Continuum path-integral notation is shorthand for the specified
finite-dimensional time slicing, finite-mode regulator, lattice regulator,
Gaussian reference measure, perturbative formal expansion, BV pushforward, or
constructive measure.  The local chapter must state which one is meant.\\
\(Z[J]\) &
Generating functional for the declared regulator or formal perturbative
scheme.  \(W[J]=\log Z[J]\) denotes the connected functional where this
logarithm is defined or formally expanded.\\
\(\Gamma[\phi]\) &
1PI effective action, defined by a Legendre transform under the local
convexity/formal-inversion assumptions.  In gauge theory, \(\Gamma\) must be
understood in the stated gauge-fixed BV or lattice framework.\\
\(S_\Lambda\), \(I_\Lambda\) &
Wilsonian action or interaction at scale \(\Lambda\).  The exact meaning
depends on the regulator and field coordinates; counterterms are part of the
Lagrangian or action specified by that regulator.\\
\(\mu\) &
Renormalization scale, not a Wilsonian cutoff unless a chapter explicitly
constructs a cutoff family with that parameter.\\
\(\Lambda\) &
Cutoff, Wilsonian scale, or lattice ultraviolet scale when supplied as part
of a regulator datum.  It is not interchangeable with \(\mu\).\\
\(\beta^i(g)\) &
Beta function in the declared coordinate chart on couplings.  Scheme changes
are coordinate changes; universal claims require the chapter to state the
invariant content.\\
\(\mathcal O_A\), \([\mathcal O_A]_\mu\) &
Local operator and renormalized local operator at scale \(\mu\).  The
operator regularization scheme and the path-integral regulator are separate
data unless the chapter identifies them.\\
\hline
\end{longtable}

\section*{Gauge-Theory and BV Symbols}
\addcontentsline{toc}{section}{Gauge-Theory and BV Symbols}

\begin{longtable}{>{\raggedright\arraybackslash}p{0.18\textwidth}
                  >{\raggedright\arraybackslash}p{0.74\textwidth}}
\hline
\textbf{Symbol} & \textbf{Meaning and scope}\\
\hline
\endhead
\(G\), \(\mathfrak g\) &
Gauge group and Lie algebra.  Global form and bundle topology are part of the
data when line operators, instantons, anomalies, or theta angles are
discussed.\\
\(A\), \(F_A\) &
Connection and curvature.  Locally \(F_A=\dd A+A\wedge A\) in
anti-Hermitian geometric conventions; component conventions are restated in
the gauge-theory chapters.\\
\(\operatorname{tr}\) &
Trace in the declared representation and normalization.  Yang--Mills
chapters use action normalization compatible with
\((4g^2)^{-1}\operatorname{tr}F_{\mu\nu}F^{\mu\nu}\) after the local
Hermitian/anti-Hermitian conversion has been stated.\\
\(c,\bar c,B\) &
Ghost, antighost, and auxiliary field in the stated gauge-fixing convention.\\
\(s\) &
BRST or BV differential.  Nilpotence may hold off shell, on shell, or after
the quantum master equation depending on the construction declared locally.\\
\(\Phi^A,\Phi_A^*\) &
BV fields and antifields.  The odd symplectic form, antibracket
\((\,\cdot\,,\,\cdot\,)\), and BV Laplacian \(\Delta\) require a chosen
finite-dimensional or regularized measure/Berezinian datum.\\
\(\mathcal W_\gamma\) &
Wilson line or Wilson loop along \(\gamma\), with representation, orientation,
and renormalization prescription stated locally.\\
\(T_\gamma\) &
't Hooft line or magnetic line datum.  Its definition requires a prescribed
singular bundle/connection behavior around \(\gamma\), not merely an
operator label.\\
\hline
\end{longtable}

\section*{CFT, Integrability, and Thermal Symbols}
\addcontentsline{toc}{section}{CFT, Integrability, and Thermal Symbols}

\begin{longtable}{>{\raggedright\arraybackslash}p{0.18\textwidth}
                  >{\raggedright\arraybackslash}p{0.74\textwidth}}
\hline
\textbf{Symbol} & \textbf{Meaning and scope}\\
\hline
\endhead
\(\Delta\), \(\ell\) &
Scaling dimension and spin of a local operator or representation.  The
letter \(\Delta\) is also used for modular operators and Laplacians in other
chapters; the local chapter disambiguates.\\
\(\mathcal O_i\) &
Local operator in a CFT or general QFT.  In CFT radial quantization, local
operators of definite scaling dimension correspond to energy eigenstates; the
linear span of local-operator states is dense in the Hilbert space under the
stated positivity/sewing hypotheses, not equal to it as a completed space.\\
\(C_{ij}{}^k\) &
OPE coefficient in a declared normalization of two-point functions and
operator basis.\\
\(\mathcal F\), \(\mathcal G\) &
Conformal block, chiral block, or form factor depending on the chapter.  The
ambient representation-theoretic decomposition is stated locally.\\
\(S_{ab}(\theta)\) &
Two-dimensional factorized scattering matrix for rapidity difference
\(\theta\), after species space, analyticity strip, unitarity, crossing, and
CDD ambiguity have been declared.\\
\(Y_a(\theta)\) &
TBA \(Y\)-function for the declared particle or string node.  Mirror-channel
and physical-channel meanings are not interchangeable.\\
\(\beta_{\rm th}\) &
Inverse temperature, often written separately from RG beta functions
\(\beta^i(g)\).\\
\(\rho_\beta\), \(\omega_\beta\) &
Thermal density matrix or KMS state.  A density matrix exists only in finite
volume or type-I settings; in AQFT the KMS state is a positive normalized
functional satisfying the KMS boundary condition.\\
\(G^R\), \(\rho_{AB}\) &
Retarded correlator and spectral function in real-time finite-temperature
chapters, with contact terms and transport limits stated locally.\\
\hline
\end{longtable}

\section*{Constructive, Lattice, and Stochastic Symbols}
\addcontentsline{toc}{section}{Constructive, Lattice, and Stochastic Symbols}

\begin{longtable}{>{\raggedright\arraybackslash}p{0.18\textwidth}
                  >{\raggedright\arraybackslash}p{0.74\textwidth}}
\hline
\textbf{Symbol} & \textbf{Meaning and scope}\\
\hline
\endhead
\(\Lambda_{\rm lat}\), \(a\) &
Finite lattice and lattice spacing.  Do not confuse a lattice region
\(\Lambda_{\rm lat}\) with a continuum cutoff scale \(\Lambda\).\\
\(U_\ell\) &
Compact lattice gauge link variable on an oriented edge \(\ell\).  Haar
measure and boundary conditions are part of the finite lattice datum.\\
\(D_{\rm W}\), \(D_{\rm ov}\) &
Wilson and overlap lattice Dirac operators when the chapter has fixed gauge
background, lattice spacing, mass parameter, and boundary conditions.\\
\(\Xi\), \(X\) &
Noise and stochastic convolution in stochastic-quantization chapters.  These
are distribution-valued random objects after a regulator/renormalization
limit has been constructed; finite smooth approximants carry regulator
subscripts.\\
\(\mathcal T\), \(\Pi\), \(\Gamma\) &
Regularity-structure model space, model realization, and reexpansion maps
when the corresponding chapter has defined them.\\
\(\mathcal K\), \(P_t\) &
Kernel or heat semigroup.  The chapter states whether it is continuum,
lattice, massive, massless, spatial, spacetime, or parabolic.\\
\(H_N\), \(P_N\) &
Finite truncation Hamiltonian and projector in Hamiltonian-truncation or
DLCQ chapters.  Continuum conclusions require an explicit extrapolation or
bound beyond finite-matrix diagonalization.\\
\(\chi^2\), \(\Sigma\) &
Correlated fit statistic and covariance matrix in numerical-evidence
chapters.  These are finite data-analysis objects; they are not continuum
proof objects unless accompanied by the stated regulator estimates and an
explicit theorem that those estimates imply.\\
\hline
\end{longtable}

\section*{Overload Warnings}
\addcontentsline{toc}{section}{Overload Warnings}

Several symbols are deliberately reused in different mathematical contexts
because the local conventions are standard.  The following overloaded
symbols require particular attention:
\[
  \Delta,\quad S,\quad Z,\quad \Gamma,\quad \rho,\quad \beta,\quad
  \Lambda,\quad \mathcal F,\quad G,\quad K .
\]
Their meanings are never inferred from typography alone.  A chapter must
declare whether, for example, \(\Delta\) is a modular operator, Laplacian,
scaling dimension, discriminant, or finite difference; whether \(S\) is an
action, S-matrix, entropy, or modular \(S\)-matrix; and whether \(Z\) is a
partition function, generating functional, wave-function residue, or central
element.  This guide records the common uses, but local declarations control
the mathematics.
